% supp_search.tex -- Appendix S3: how the search works; two observations made with the search program alone
% (Sections 6.2 and 6.5 of the first revision)
\section{The search, and two observations about the search space}\label{sec:supp-search}

This appendix gives the details of the search that produced the certificate, which \cref{ssec:search}
describes in outline, and of the two observations that are summarized in \cref{rem:observations}. The
search is a heuristic device and is not part of the proof: only its output is used, and the output is
verified from scratch by programs that share no code with it (\cref{sec:verification}). We describe the
search because a reader who wants to reproduce or extend the computation will need to know how a tree of
this size was found.

\subsection{How the search works}\label{ssec:supp-search-details}

The search program is \fn{src/search.cpp}.

\subsubsection*{The basic recursion}
The search is a depth-first evaluation of the definition of \WON{} (\cref{def:won}). A P-state is an AND
node: it is proved when it satisfies $\LEAF(B,F,t)$, or when $t<\Rmax$, $\LOST$ is false, and its child~(a)
and all its non-\VOID{} threat children~(c) and danger children~(d) are proved. A Q-state is an OR node over
the table move $m$; for a given $m$ both the child~(i) (the cell is empty and Breaker takes it) and, if
unrevealed stones exist and the child is not \VOID, the child~(ii) (the cell turns out to be occupied and is
revealed) have to be proved. Every state is replaced by its canonical form (\cref{lem:canon}) when it is
entered, and results are stored in a transposition table keyed by the canonical state, so that each
symmetry class is treated once. At a P-state the leaf test comes first; it is preceded by a cheap
screen (the total surplus of all live lines is compared with $\theta=t-|F|$) and is performed with the
max-flow procedure of \cref{alg:leaf}.

\subsubsection*{Candidate moves at Q-states}
For the Q-state $\Qst(B,F,t)$ the candidate table moves are ordered as follows. Let
$\Busy$ denote the set of busy cells (\cref{def:busy}). Call a set $B'\subseteq\Busy$ of size $k\in\{8,9\}$
\emph{good} if $\LEAF(B',\emptyset,k)$ holds; the good sets are known in advance (\cref{ssec:leaftable}).
\begin{enumerate}[label=(\arabic*),leftmargin=2.2em]
\item Busy cells $m\notin B\cup F$ are sorted by: first, whether $\Pst(B\cup\{m\},F,t)$ is immediately a
  leaf; second, whether $m$ is the antipode of a stone of $B$; third, the number of good $8$-sets that
  contain $(B\cap\Busy)\cup\{m\}$ and avoid $F$, then the same number for good $9$-sets; fourth, the slack
  of $B\cup\{m\}$ (\cref{def:slack}); ties are broken by the cell id.
\item A second list contains cells that need not be busy. When the leaf test for $(B,F,t)$ fails, the
  max-flow procedure returns a witness: a set $S$ of live lines that violates Hall's condition, or an
  $F$-connected set $S$ with $\sur_F(S)<\theta$. The cells outside $B\cup F$ are ranked by the number of
  lines of $S$ through them (then by the number of live lines through them), and the first six are kept.
\item The final order is: immediate leaves first, then the two lists interleaved.
\end{enumerate}
Only the first $10$ candidates of this order are considered at all, and at most $4$ of them are examined in
full (see below). These limits make the search incomplete: a failure of the search does not show that a
state is not \WON.

\subsubsection*{The relaxed pre-pass}
The decisive device for speed is a \emph{relaxed game}: the same game without $\LOST$ and without the
forced children (c) and (d), i.e.\ with the single child (a) at P-states. With the same candidate lists, a
state that is \WON{} is also \WON{} in the relaxed game, by induction on the rank of the state
(\cref{lem:rank}): the relaxed game imposes a subset of the requirements. The search therefore evaluates
the relaxed game
first (with its own entries in the transposition table) and uses it only to fail fast: a P-state whose
relaxed version fails is abandoned, and a candidate move whose children (i) or (ii) fail in the relaxed
game is skipped; at most $4$ candidates that survive this filter are tried in the real game. Inside the
relaxed game the program prunes further with a bound that we state without proof: a stone on a cell with
$d$ live lines can raise the difference
between the total surplus and $\theta$ by at most $2d-2$, so a deficit that the best $\Rmax-t$ remaining
stones cannot close is treated as hopeless. Relaxed results are never written out and can only produce
failures.

\subsubsection*{Other details}
At an expanded P-state the forced children are searched before the child (a). The forced children of
P-states on the main trunk are distributed over threads; the transposition table is shared and lock-free
and stores $112$ hash bits per state. When the root has been proved, the proof tree is walked once more
from the root and the Q-states with their moves and the leaves that occur in it are written to the two
files. With the option \fn{--selfcheck} the program then reads the files back and re-expands the tree with
a deliberately naive child generator.

\subsubsection*{The run}
The certificate was produced by
\begin{verbatim}
bin/search --rmax 10 --maxcand 4 --relcand 10 --nonbusy --threads 4 \
           --out results/final --selfcheck
\end{verbatim}
in $57.9$ seconds of wall-clock time on a four-core machine (\cref{sec:repro}), with a peak memory of
$601$\,MB, most of which is the fixed-size transposition table. The search expanded $49{,}397$ P-states
and $96{,}487$ Q-states of the real game and $48{,}937$ and $97{,}478$ of the relaxed game, made
$655{,}010$ leaf tests ($25.7$ CPU-seconds) and created $930{,}120$ transposition-table entries
(\fn{results/final/stats.json}). Since $46{,}412$ of the $49{,}397$ expanded P-states belong to the final
tree, the search does very little backtracking with this move ordering. What is known about the
dependence of the output on the number of threads is stated in \cref{ssec:search}.

\subsection{Computed observations about the search space}\label{ssec:supp-observations}

The following two observations were obtained with the search program alone. They are \emph{not} verified
by an independent program, they are not used anywhere in the proof, and we record them only because they
explain two features of the certificate.

\subsubsection*{Busy cells alone seem not to suffice for $\Rmax\le 10$ (single program)}
It is natural to restrict Breaker's table moves to the $24$ busy cells, each of which lies on seven
lines. According to the search program, the root is not \WON{} with this restriction for $\Rmax=10$, and
hence for no smaller $\Rmax$ (a
valid certificate for $\Rmax$ is valid for every larger value): already the relaxed game of
\cref{ssec:supp-search-details}, in which Maker never forces, fails. For this run the search tried \emph{all} busy cells
outside $B\cup F$ at every Q-state, so that it is exhaustive for the relaxed game with busy table moves:
\begin{verbatim}
bin/search --rmax 10 --maxcand 4 --relcand 24      (no --nonbusy)
\end{verbatim}
reports failure after visiting $12{,}437$ relaxed P-states and $14{,}196$ relaxed Q-states ($14$ to $22$
seconds in three runs, one of them made during the internal refereeing of this paper with a binary built
from the archive; all runs gave the same counts). The conclusion that busy cells do not suffice rests on three things: the
correctness of the search program in this mode, which no second program has checked; the fact that a
state that is \WON{} is also \WON{} in the relaxed game (\cref{ssec:supp-search-details}); and the pruning bound of the
relaxed game, which was applied $48$ times in this run and which is not proved in this paper.

The obstruction is easy to describe. In the failing branch Maker's revealed stones are the four corners
of a face, for instance $F=\{(5,0,0),(5,0,5),(5,5,0),(5,5,5)\}$ in the face $x=5$. The $14$ lines that lie
inside a face ($6$ rows, $6$ columns and $2$ diagonals) can only be killed by a stone of Breaker inside
that face, and the only busy cells of a face are its four corners. As long as Breaker plays busy cells,
the set
$S$ of these $14$ lines therefore stays alive; it is $F$-connected, and
\[
\sur_F(S)=|\N(S)\setminus F|-2|S| = (36-4)-28 = 4 .
\]
By the definition of $\LEAF$ (\cref{def:leaf}), a leaf with these revealed stones needs
$\sur_F(S)\ge\theta=t-4$, i.e.\ $t\le 8$; and by round $8$ Breaker has not reached a leaf in this branch
either. (In the final position of the failing branch, after round $10$, the max-flow procedure indeed
returns exactly these $14$ lines as its witness, with surplus $4<\theta=6$.) This is why the search was run
with the option \fn{--nonbusy} and why $434$ table moves of the certificate are non-busy cells. For
$\Rmax=11$ the search with busy cells also failed, but with only $8$ candidates per state, which is not
exhaustive. Whether busy table moves suffice for some larger $\Rmax$ is open.

\subsubsection*{$\Rmax=9$ was not achieved}
With non-busy moves allowed and $\Rmax=9$ the search fails, both with the candidate limits of the
certified run and with up to $6$ candidates per state (options \texttt{-{}-rmax 9 -{}-maxcand 6 -{}-relcand 16 -{}-nbextra 10 -{}-nonbusy};
$128$ seconds). The failing branch is a double threat: Maker makes a first threat that is answered in round
$5$ and a second one, through a stone of the first line, that is answered in round $9$; the resulting P-state
of round $9$ is not a leaf. Since the candidate lists are truncated, this does not show that the root is not
\WON{} for $\Rmax=9$. Whether $\Rmax=9$ suffices is open; the theorem does not depend on it.
